\begin{theorem}[6.8.6]
\label{IV.2.s6.8.6-fr}
Sean $X$, $Y$ dos preesquemas localmente noetherianos, $f:X\to Y$ un
morfismo localmente de tipo finito, $x$ un punto de $X$, $y=f(x)$. Las
propiedades siguientes son equivalentes:
\begin{enumerate}
\item[\textup{a)}] $f$ es suave en el punto $x$.
\item[\textup{b)}] $f$ es regular en el punto $x$.
\oldpage[IV-2]{153}
\item[\textup{c)}] $\mathcal O_x$ es una $\mathcal O_y$-álgebra formalmente suave
\textup{(\textbf 0, 19.3.1)} para las topologías preádicas en $\mathcal O_x$ y $\mathcal O_y$.
\item[\textup{c')}] $\mathcal O_x$ es una $\mathcal O_y$-álgebra formalmente suave
\textup{(\textbf 0, 19.3.1)} para las topologías discretas en $\mathcal O_x$ y $\mathcal O_y$.
\end{enumerate}
\end{theorem}

\begin{proof}
La equivalencia de \emph{a)} y \emph{b)} resulta de las definiciones (6.8.1) y del hecho de que
para los preesquemas localmente noetherianos, los morfismos localmente de tipo
finito son localmente de presentación finita (1.4.2).

En segundo lugar, para que
$\mathcal O_x$ sea una $\mathcal O_y$-álgebra formalmente suave para las
topologías preádicas, es necesario y suficiente que $\mathcal O_x$ sea un
$\mathcal O_y$-módulo plano y que $\mathcal O_x\otimes_{\mathcal O_y}\boldsymbol{k}(y)$
sea una $\boldsymbol{k}(y)$-álgebra formalmente suave para su topología preádica
(\textbf 0, 19.7.1); pero para que $\mathcal O_x\otimes_{\mathcal O_y}\boldsymbol{k}(y)$ sea una
$\boldsymbol{k}(y)$-álgebra formalmente suave para su topología preádica,
es necesario y suficiente que sea una $\boldsymbol{k}(y)$-álgebra geométricamente regular
(\textbf 0, 19.6.6); esto demuestra la equivalencia de \emph{b)} y \emph{c)}. Por último, para probar
la equivalencia de \emph{c)} y \emph{c')}, es posible limitarse al caso en que $Y=\operatorname{Spec}(A)$
y $X=\operatorname{Spec}(C)$ son afines, donde $A$ es noetheriano y $C$ una
$A$-álgebra de tipo finito, que por tanto se puede escribir en la forma
$A[T_1,\ldots,T_n]/\mathfrak J$. Como aquí $\mathfrak J/\mathfrak J^2$ es un
$C$-módulo de presentación finita, la equivalencia de \emph{c)} y \emph{c')} resulta de
(\textbf 0, 22.6.4) aplicado a $A$, $B=A[T_1,\ldots,T_n]$ y $C=B/\mathfrak J$.
\end{proof}

\begin{corollary}[6.8.7]
\label{IV.2.s6.8.7-fr}
Sean $X$, $Y$ dos preesquemas localmente noetherianos, $f:X\to Y$ un
morfismo localmente de tipo finito. Entonces el conjunto de puntos $x\in X$ en los que
$f$ es suave (\textup{o} regular) es abierto en $X$.
\end{corollary}

\begin{proof}
En efecto, de (\textbf 0, 22.6.5) resulta que el conjunto de los $x\in X$ que satisfacen la
condición \emph{c')} de (6.8.6) es abierto en $X$, y se concluye por (6.8.6).
\end{proof}

\begin{remark}[6.8.8]
\label{IV.2.s6.8.8-fr}
En (17.5.1) mostraremos que la equivalencia de \emph{b)} y \emph{c')} en (6.8.6),
así como el corolario (6.8.7), siguen siendo válidos sin la hipótesis noetheriana
en $X$ e $Y$, siempre que se consideren únicamente morfismos \emph{localmente de
presentación finita}.
\end{remark}
